% Background

\subsection{Redistricting as a Decadal Process}

The U.S. Constitution mandates that congressional seats be reapportioned among states following each decennial census, with districts redrawn to reflect new population counts. This creates a natural 10-year cycle that has operated consistently since 1790. The three most recent redistricting cycles—2000, 2010, and 2020—occurred during distinct political and technological contexts:

\begin{itemize}
    \item \textbf{2000}: Pre-REDMAP era, traditional redistricting methods
    \item \textbf{2010}: Post-REDMAP era, sophisticated optimization enabled by modern GIS tools
    \item \textbf{2020}: Post-commission era, with several states adopting independent redistricting commissions
\end{itemize}

\subsection{Prior Longitudinal Work}

Several studies have examined redistricting trends over time. \citet{altman2011gerrymandering} documented increasing partisan bias from 1972--2010 using electoral outcomes. \citet{henderson2018measuring} analyzed compactness measures across multiple redistricting cycles but used different analytical approaches for different years. \citet{mcghee2014redistricting} examined efficiency gaps across decades, focusing on partisan metrics. Our approach complements this work by examining geometric properties (compactness, boundary stability) rather than partisan outcomes.

\textbf{Gap}: No prior work has applied a consistent algorithmic methodology across multiple census decades, making it impossible to separate genuine temporal trends from methodological variations.

\subsection{Algorithmic Redistricting}

Algorithmic redistricting uses computational optimization to generate districts based on explicit criteria (population balance, compactness, contiguity) without partisan data. The recursive bisection approach \citep{karypis1998metis} treats census tracts as a graph and recursively partitions it to minimize edge cuts (maximizing compactness) while balancing population.

\textbf{Key advantage for longitudinal analysis}: By applying the identical algorithm with identical parameters across 2000, 2010, and 2020 data, we isolate temporal effects from methodological changes.

\subsection{Research Program Context}

This paper extends our prior work on recursive bisection redistricting:
\begin{itemize}
    \item \textbf{Paper 01}: Established recursive bisection as the core methodology
    \item \textbf{Paper 05}: Developed cross-census validation techniques
    \item \textbf{Paper 07}: Examined temporal stability within single census years
\end{itemize}

This paper represents the first \textit{inter-census} longitudinal analysis across 20 years.
